
In this chapter, we performed a comparative evaluation of eight unsupervised ML models applied to the detection of users' QoE degradation in CG sessions. Our evaluation showed that the F1-score metric, which is widely used for model evaluation, has limitations and should be combined with the MCC score for more accurate model evaluation. Moreover, existing window-based approaches, used to cover sequences of events including anomalies, lead to erroneous conclusions regarding model performance. To address these limitations, we proposed the WAD approach to allow for a fair and better assessment. The WAD approach offers the possibility of parameter calibration to detect more or less severe anomalies. As practical considerations in the use of ML models in networking, we also showed that data contamination has a considerable impact on unsupervised ML models, and revealed their disparity with respect to their computational performance. Our study has shown that the use of models such as those included in our evaluation in an industrial context requires further investigation of their applicability and calibration. Highly performing state-of-the art ML models have not been necessarily designed with industrial considerations in mind such as robustness, energy consumption, explainability and likely others. Many of these considerations can be conflicting with commonly used performance evaluation metrics. 
In summary, this chapter emphasizes the importance of employing a consistent methodology and appropriate metrics when evaluating ML models. In particular, we found no one-size-fits-all solution as some solutions may be preferred to others depending on the requirements of the operational environment under consideration. These insights will serve as the foundation for subsequent chapters, where we further propose novel anomaly detection and root-cause diagnosis methods in the context of low-latency applications.



% As future work, we plan to investigate and develop anomaly detection models capable to achieve the right trade-offs between these seemingly conflicting goals when approaching the detection of abnormal network events and their impact on user QoE particularly in the case of demanding applications  such as low-latency applications. 
% To address these objectives, we will specifically focus on QoE degradation detection based on network packets collected at the network edge. We will collect these packets in various experimental and in-production scenarios, which will enable us to build more robust and generalizable ML models. Furthermore, we will conduct root-cause analysis for the identified abnormal network events, which is critical to automatically initiate mitigation operations. Our ultimate aim is to advance the state-of-the-art in anomaly detection techniques for efficient troubleshooting of low-latency applications and for minimizing the impact of abnormal events on user QoE.